\documentclass[10pt,a4paper]{amsart}
\usepackage[margin=19mm]{geometry}
\usepackage{amsmath,amssymb,mathtools}
\usepackage[scr=rsfs,cal=euler]{mathalfa}
\usepackage{xfrac}
\usepackage[unicode,hidelinks,bookmarks=false]{hyperref}
\newcommand{\ie}{i.e.\ }
\newcommand{\cf}{cf.\ }
\newcommand{\resp}{respectively\ }
\newcommand{\Subsection}{\subsection}
\providecommand{\nobreakdash}{\nobreak\hbox{-}}
\setlength{\emergencystretch}{2em}
\multlinegap0pt
\usepackage{bm}
\usepackage{bbm}
\usepackage{dsfont}
\usepackage{xspace}
\usepackage{cancel}
\usepackage{mathtools}
\usepackage{amsthm}
\makeatletter
\@ifundefined{proposition}{\newtheorem{proposition}{Proposition}}{}
\makeatother
\title[Extending Hridaya Kolam to Multiple Loops: A Study of Non-Co]{Extending Hridaya Kolam to Multiple Loops: A Study of Non-Coprime Dot--Arm Structures}
\author{Atanu Manna, Suvra Kanti Chakraborty}
\date{Source: arXiv:2510.18907v1 (CC BY 4.0)}
\begin{document}
\maketitle

\noindent\emph{Source and licence.} Extending Hridaya Kolam to Multiple Loops: A Study of Non-Coprime Dot--Arm Structures. Version arXiv:2510.18907v1 (CC BY 4.0), distributed under the Creative Commons Attribution 4.0 International licence, as recorded in the arXiv licence metadata for this exact version and shown on the abstract page https://arxiv.org/abs/2510.18907v1. Full rights notice: licenses/arxiv-2510.18907-CC-BY-4.0.txt.\par\smallskip

\noindent\textit{Editorial scope.} Counted unit: the Proposition whose source label is \texttt{None} in the article (source0.tex lines 209--218, source label \texttt{None}), delivered together with the article's own complete original proof of it (source0.tex lines 220--270, one \texttt{proof} environment of 51 lines). No mathematics is written, repaired or re-derived: statement and proof are the article's own text line for line. Required context retained uncounted: none. Every definition, display and label of those lines is retained unchanged. Omitted from this package and not redistributed: 1--208, 219--219, 271--811. Nothing inside a retained excerpt is omitted or altered: every retained body is delivered exactly as the article's lines are written, so no omission receipt of any kind accompanies this package. Citations: the retained excerpts carry no citation command at all, so no bibliography entry of the article is transcribed and no reference is redistributed. Reference adaptation: none was needed and none was made; every reference inside the delivered unit resolves inside it, so no unresolved reference or citation is produced. Printed slips kept literal: no printed word, symbol, hypothesis or number of the counted statement or of its proof was altered. Layout and class: the article's own class is \texttt{article} with packages including algorithm, algpseudocode, inputenc, fontenc, lmodern, enumitem, amsmath, amssymb, amsthm, authblk, float, array, booktabs, graphicx; the wrapper is the lane's pinned \texttt{amsart} editorial wrapper at 10pt on A4, which restates the theorem-like environments the retained text needs and adds the article's own macro definitions that the retained text needs (none), with extra wrapper packages (bm, bbm, dsfont, xspace, cancel, mathtools). The delivered numbering is the unit's own. Assets: no retained span contains a figure environment, a graphics-inclusion command, a PDF-page inclusion, a table or a TikZ picture; no image file is copied into this package, the package declares no asset dependency and no image file is redistributed with it. Licence: version arXiv:2510.18907v1 of this article is distributed under the Creative Commons Attribution 4.0 International licence, as recorded in the arXiv licence metadata for this exact version and shown on the abstract page https://arxiv.org/abs/2510.18907v1. A full rights notice is delivered at licenses/arxiv-2510.18907-CC-BY-4.0.txt. Characters: every retained excerpt is pure ASCII, so the delivered engine, XeLaTeX with its default Latin Modern text font, reports no missing character.
\begin{proposition}
Let $m,n$ be positive integers and let $d=\gcd(m,n)>1$. 
For each residue $r\in\{0,1,\dots,d-1\}$ define the sequence
\[
a^{(r)}_k \equiv (k n + r) \mod m, \qquad k=0,1,2,\dots
\]
where residues modulo $m$ are represented by $\{0,1,\dots,m-1\}$. 
Then the $m$ integers $0,1,\dots,m-1$ are divided into exactly $d$ disjoint repeating sequences of this form. 
Hence, the number of disjoint loops produced by the construction equals $d=\gcd(m,n)$.
\end{proposition}

\begin{proof}
Let $d=\gcd(m,n)$. 
Then we may write $m=d\,m_1$ and $n=d\,n_1$ where $\gcd(m_1,n_1)=1$.

\noindent Consider the first sequence corresponding to $r=0$:
\[
a^{(0)}_k \equiv (k n) \mod m = (k d n_1) \mod{(d m_1)}.
\]
Two terms $a^{(0)}_k$ and $a^{(0)}_{k'}$ are equal modulo $m$ if and only if 
\[
m \mid (k-k')n \quad \iff\quad d m_1 \mid (k-k')d n_1 
\quad \iff\quad m_1 \mid (k-k')n_1.
\]
Since $\gcd(m_1,n_1)=1$, this happens exactly when $m_1\mid (k-k')$. 
Therefore the sequence repeats after every $m_1$ steps, i.e.
\[
a^{(0)}_{k+m_1} \equiv a^{(0)}_k \mod m.
\]

Hence, the sequence $a^{(0)}_k$ contains exactly $m_1 = \dfrac{m}{d}$ distinct values before repeating. 
Similarly, generalizing this for any $r$, it can be shown that the sequence $a^{(r)}_k$ also contains exactly $\dfrac{m}{d}$ distinct elements before repetition.\\


\noindent Now observe that for any $r\in\{0,1,\dots,d-1\}$,
\[
a^{(r)}_k \equiv (k n + r) \mod m = (k d n_1 + r) \mod{d m_1}.
\]
\noindent
So we can write
\[
a^{(r)}_k = k d n_1 + r + t (d m_1)
\]
for some integer \(t\).\\


\noindent
Looking at the right-hand side, we note that \(k d n_1\) is divisible by \(d \;\Rightarrow\; k d n_1 \equiv 0 \pmod{d}\) and \(t (d m_1)\) is also divisible by \(d \;\Rightarrow\; t d m_1 \equiv 0 \pmod{d}\) 

\noindent
Therefore,
\[
a^{(r)}_k \equiv r \pmod{d}.
\]

\noindent Consequently, two sequences $a^{(r_1)}_k$ and $a^{(r_2)}_k$ can never have a common term unless $r_1\equiv r_2 \mod d$. 
Since we have taken $r=0,1,\dots,d-1$, the $d$ sequences are pairwise disjoint.

\noindent Each sequence contains exactly $\dfrac{m}{d}$ distinct elements, and all $d$ sequences together cover all $m$ integers $0,1,\dots,m-1$. 
Therefore, the total number of disjoint sequences (or loops) is exactly $d=\gcd(m,n)$.
 
\end{proof}

\end{document}
